%
% 第2章：激活函数

\clearpage
\cleardoublepage

\chapter{激活函数}

\section{引言}

激活函数是应用于神经网络中每个神经元（或层）输出的数学函数。它们向网络引入\textbf{非线性}，使其能够学习线性函数无法捕获的复杂模式。

激活函数的选择显著影响：
\begin{itemize}
    \item 网络学习复杂模式的能力
    \item 训练稳定性和收敛速度
    \item 反向传播中的梯度流动
    \item 神经元的输出范围
    \item 计算效率
\end{itemize}

\section{Sigmoid函数族}

Sigmoid函数将输入压缩到范围 $(0, 1)$。

\subsection{Sigmoid（逻辑函数）}

\begin{equation}
\sigma(x) = \frac{1}{1 + e^{-x}} = \frac{e^x}{e^x + 1}
\end{equation}

属性：
\begin{itemize}
    \item \textbf{输出范围：} $(0, 1)$
    \item \textbf{导数：} $\sigma'(x) = \sigma(x)(1 - \sigma(x))$
    \item \textbf{可解释为概率}
\end{itemize}

\begin{advantages}
    \item 平滑梯度
    \item 输出可解释为概率
    \item 历史上广泛使用
\end{advantages}

\begin{disadvantages}
    \item \textbf{梯度消失：} 对于大的 $|x|$，$\sigma'(x) \approx 0$
    \item \textbf{非零中心：} 输出总是正的
    \item \textbf{饱和：} 在0和1处饱和
\end{disadvantages}

\subsection{双曲正切（Tanh）}

\begin{equation}
\tanh(x) = \frac{e^x - e^{-x}}{e^x + e^{-x}} = 2\sigma(2x) - 1
\end{equation}

属性：
\begin{itemize}
    \item \textbf{输出范围：} $(-1, 1)$
    \item \textbf{导数：} $\tanh'(x) = 1 - \tanh^2(x)$
    \item \textbf{零中心输出}
\end{itemize}

比Sigmoid更强的梯度，通常是隐藏层的更好选择。

\section{ReLU函数族}

ReLU函数族因计算效率高和训练深度网络的有效性而成为现代深度学习的主流选择。

\subsection{ReLU}

\begin{equation}
\text{ReLU}(x) = \max(0, x) = 
\begin{cases}
x & \text{if } x > 0 \\
0 & \text{if } x \leq 0
\end{cases}
\end{equation}

\begin{properties}[ReLU属性]
\begin{itemize}
    \item \textbf{输出范围：} $[0, \infty)$
    \item \textbf{导数：} $\text{ReLU}'(x) = \begin{cases} 1 & x > 0 \\ 0 & x \leq 0 \end{cases}$
    \item \textbf{计算高效}
    \item \textbf{稀疏激活}
\end{itemize}
\end{properties}

\begin{advantages}
    \item \textbf{计算非常快}
    \item \textbf{减少正激活的梯度消失问题}
    \item \textbf{稀疏表示}
    \item \textbf{深度网络的出色经验性能}
\end{advantages}

\begin{disadvantages}
    \item \textbf{Dying ReLU问题：} 神经元可能"死亡"
    \item \textbf{非零中心}
    \item \textbf{输出无界}
\end{disadvantages}

\begin{note}[Dying ReLU问题]
如果神经元的权重更新使得所有输入都满足 $\mathbf{w}^\top \mathbf{x} + b \leq 0$，神经元将永远输出零。这被称为"死亡ReLU"问题。
\end{note}

\subsection{Leaky ReLU}

\begin{equation}
\text{LeakyReLU}(x) = 
\begin{cases}
x & \text{if } x > 0 \\
\alpha x & \text{if } x \leq 0
\end{cases}
= \max(0, x) + \alpha \min(0, x)
\end{equation}
其中 $\alpha$ 是一个小正数（通常 $\alpha = 0.01$）。

解决了Dying ReLU问题，同时保持ReLU的优点。

\subsection{GELU（高斯误差线性单元）}

\begin{equation}
\text{GELU}(x) = x \cdot \Phi(x)
\end{equation}
其中 $\Phi(x)$ 是标准正态分布的CDF。

GELU用于BERT、GPT和许多Transformer架构。

\begin{properties}[GELU属性]
\begin{itemize}
    \item \textbf{随机正则化器}
    \item \textbf{所有区域非线性}
    \item \textbf{接近零中心}
\end{itemize}
\end{properties}

\section{Softmax函数}

\begin{equation}
\text{softmax}(x_i) = \frac{e^{x_i}}{\sum_{j=1}^{K} e^{x_j}} \quad \text{for } i = 1, \ldots, K
\end{equation}

Softmax用于多类分类输出层。

属性：
\begin{itemize}
    \item \textbf{输出范围：} 每个输出在 $(0, 1)$
    \item \textbf{和约束：} $\sum_i \text{softmax}(x_i) = 1$
    \item \textbf{概率解释}
\end{itemize}

\section{选择激活函数}

\begin{table}[h]
    \centering
    \caption{激活函数选择指南}
    \begin{tabular}{L{3cm} L{5cm} L{4cm}}
        \toprule
        \textbf{层类型} & \textbf{推荐函数} & \textbf{备注} \\
        \midrule
        \textbf{隐藏层（通用）} & ReLU, GELU & ReLU是默认选择 \\
        \textbf{隐藏层（深度）} & ReLU, GELU, ELU & 避免深度网络中的sigmoid/tanh \\
        \textbf{隐藏层（RNN）} & Tanh & 通常是RNN的默认值 \\
        \textbf{隐藏层（Transformer）} & GELU & BERT、GPT等使用 \\
        \textbf{输出（二分类）} & Sigmoid & 概率输出 \\
        \textbf{输出（多分类）} & Softmax & 概率分布 \\
        \textbf{输出（回归）} & Identity / 无 & 无界输出 \\
        \bottomrule
    \end{tabular}
\end{table}

\section{本章小结}

\begin{enumerate}
    \item \textbf{激活函数引入非线性}，使网络能够学习复杂模式。
    
    \item \textbf{Sigmoid和Tanh} 将输入压缩到有界范围，但存在梯度消失问题。
    
    \item \textbf{ReLU} 计算效率高，减少梯度消失，但可能存在死亡神经元问题。
    
    \item \textbf{Leaky ReLU、PReLU和ELU} 解决了Dying ReLU问题。
    
    \item \textbf{GELU} 在Transformer架构中变得流行。
    
    \item \textbf{Softmax} 用于多类分类输出层。
    
    \item \textbf{选择取决于} 任务、架构和训练稳定性。
\end{enumerate}
